\documentclass{article}
\usepackage{booktabs}
\usepackage{array}
\usepackage{caption}
\usepackage{geometry}
\usepackage{placeins}
\usepackage{graphicx}
\usepackage{grffile}
\geometry{landscape, margin=1in}

\begin{document}

\begin{figure}[htbp]
  \centering
  \includegraphics[width=0.9\textwidth]{\detokenize{../Results/Fig_SharpeRatio_OOSvsIS_ByCohort.pdf}}
  \caption{Out-of-sample to in-sample Sharpe ratio by sample-end decade cohort: published vs data-mined. For each published signal and its matched data-mined portfolio, we compute the annualized Sharpe ratio separately over the in-sample and post-sample periods. The Sharpe ratio is invariant to the return scaling used elsewhere in the paper, so results are identical whether computed from raw or scaled returns. The box plot shows the distribution of the OOS/IS Sharpe ratio across signals within each cohort. A ratio of 1 (dashed line) indicates no deterioration; values below 1 indicate that risk-adjusted performance declined out of sample. The data-mined Sharpe ratio is computed from the time series of cross-sectional average returns of matched Compustat accounting strategies with $|t|>2.0$ in the original sample period. Following Campbell (2024), signals with fewer than 5 years of post-sample data are excluded, cohorts with fewer than 10 signals are dropped, and both periods require at least 24 months.}
\end{figure}

\begin{figure}[htbp]
  \centering
  \includegraphics[width=0.9\textwidth]{\detokenize{../Results/Fig_SharpeRatio_OOSvsIS_ByTheory.pdf}}
  \caption{Out-of-sample to in-sample Sharpe ratio by theory category: published vs data-mined. For each published signal and its matched data-mined portfolio, we compute the annualized Sharpe ratio separately over the in-sample and post-sample periods. Signals are grouped by their theoretical motivation: ``risk'' (compensation for systematic risk), ``mispricing'' (behavioral or limits-to-arbitrage explanations), or ``agnostic'' (no specific theoretical motivation). A ratio of 1 (dashed line) indicates no deterioration. The data-mined Sharpe ratio is computed from the time series of cross-sectional average returns of matched Compustat accounting strategies with $|t|>2.0$ in the original sample period. Following Campbell (2024), signals with fewer than 5 years of post-sample data are excluded, and cohorts with fewer than 10 signals are dropped.}
\end{figure}

\begin{figure}[htbp]
  \centering
  \includegraphics[width=0.9\textwidth]{\detokenize{../Results/Fig_SharpeRatio_SR2Weighted_PubVsDM_OOS.pdf}}
  \caption{Mean-variance optimal portfolio: published vs data-mined returns in calendar time (at least five years post-sample). Following Campbell (2024), instead of scaling each signal's returns to a common in-sample mean and equal-weighting (which implicitly weights inversely by squared Sharpe ratio), we weight each signal by its in-sample mean-variance optimal weight $\mu_{IS}/\sigma_{IS}^2$. Published and data-mined returns are each weighted by their own in-sample $\mu/\sigma^2$. For each calendar month, we compute the weighted average of long-short returns across all published signals (and separately, their matched data-mined strategies) that are at least five years past their original-sample end date. The plot shows a trailing 60-month (5-year) rolling mean. Signals with fewer than 5 years of post-sample data are excluded, and signals in decade cohorts with fewer than 10 members are dropped.}
\end{figure}

\begin{figure}[htbp]
  \centering
  \includegraphics[width=0.9\textwidth]{\detokenize{../Results/Fig_SharpeRatio_SR2Weighted_ByCohort_OOS.pdf}}
  \caption{Mean-variance optimal published signal returns by sample-end decade cohort in calendar time (at least five years post-sample). Each published signal is weighted by its in-sample $\mu_{IS}/\sigma_{IS}^2$. For each cohort and calendar month, we compute the weighted average of long-short returns across signals in that cohort that are at least five years past their original-sample end date. Faded colored lines show the trailing 5-year mean return for each cohort; the bold black line shows the weighted average across all signals. Following Campbell (2024), signals with fewer than 5 years of post-sample data are excluded, and cohorts with fewer than 10 signals are dropped.}
\end{figure}

\begin{figure}[htbp]
  \centering
  \includegraphics[width=0.9\textwidth]{\detokenize{../Results/Fig_SharpeRatio_SR2Weighted_DM_ByCohort_OOS.pdf}}
  \caption{Mean-variance optimal data-mined signal returns by sample-end decade cohort in calendar time (at least five years post-sample). Each data-mined portfolio is weighted by its own in-sample $\mu_{IS}/\sigma_{IS}^2$. For each cohort and calendar month, we compute the weighted average of long-short returns from the matched data-mined Compustat accounting strategies with $|t|>2.0$ in the original sample period, using only observations at least five years past each signal's original-sample end date. Faded colored lines show the trailing 5-year mean return for each cohort; the bold black line shows the weighted average across all signals. Following Campbell (2024), signals with fewer than 5 years of post-sample data are excluded, and cohorts with fewer than 10 signals are dropped.}
\end{figure}

\FloatBarrier

\end{document}
